\documentclass[12pt]{article}

% Packages
\usepackage[a4paper, margin=1in]{geometry}
\usepackage{setspace}
\usepackage{amsmath, amssymb}
\usepackage{titlesec}
\usepackage{enumitem}
\usepackage{hyperref}

% Formatting
\onehalfspacing
\setlist[itemize]{noitemsep}
\setlist[enumerate]{noitemsep}

\titleformat{\section}{\large\bfseries}{\thesection.}{0.5em}{}
\titleformat{\subsection}{\normalsize\bfseries}{\thesubsection}{0.5em}{}

\begin{document}

\begin{center}
{\Large \textbf{CSE400 -- Fundamentals of Probability in Computing}}\\[0.2cm]
{\large Lecture 19: N-Variate Joint PDF / PMF / CDF}\\[0.3cm]

\textbf{Topic:} Joint Distribution of N Random Variables\\
\end{center}

\hrule
\vspace{0.4cm}

\section{Joint Distribution of N Random Variables}

\subsection{Definition}

Let $X_1, X_2, \dots, X_n$ be $n$ random variables defined on a common probability space.

The joint cumulative distribution function (CDF) is defined as:

\[
F_{X_1,\dots,X_n}(x_1,\dots,x_n) =
P(X_1 \le x_1, X_2 \le x_2, \dots, X_n \le x_n)
\]

\subsubsection*{Step-by-Step Reasoning}

\begin{enumerate}
\item Consider each random variable individually.
\item Form the joint event:
\[
(X_1 \le x_1) \cap (X_2 \le x_2) \cap \cdots \cap (X_n \le x_n)
\]
\item Compute the probability of this intersection.
\item This defines the joint CDF.
\end{enumerate}

\section{Example 1}

\subsection{Reasoning Procedure}

\begin{enumerate}
\item Identify all random variables.
\item Construct the joint outcome space.
\item Assign PMF (discrete) or PDF (continuous).
\item Verify normalization:
\end{enumerate}

\textbf{Discrete case:}
\[
\sum_{x_1} \cdots \sum_{x_n} p(x_1,\dots,x_n) = 1
\]

\textbf{Continuous case:}
\[
\int f(x_1,\dots,x_n)\, dx_1 \cdots dx_n = 1
\]

\section{Application: Large MIMO System}

\subsection{Interpretation}

Multiple antennas correspond to multiple random variables, forming an $n$-dimensional random vector.

\subsection{Step-by-Step Reasoning}

\begin{itemize}
\item Let each antenna output be $X_i$.
\item Form the vector $(X_1, X_2, \dots, X_n)$.
\item Describe system behavior using joint distribution.
\item Capture dependencies among variables.
\end{itemize}

\section{Joint PMF, CDF, PDF}

\subsection{Joint PMF (Discrete Case)}

\[
p(x_1,\dots,x_n) =
P(X_1 = x_1, \dots, X_n = x_n)
\]

\subsubsection*{Reasoning}

\begin{itemize}
\item Each tuple represents an outcome.
\item Assign probability to each tuple.
\item Ensure normalization:
\end{itemize}

\[
\sum_{x_1} \cdots \sum_{x_n} p(x_1,\dots,x_n) = 1
\]

\subsection{Joint CDF}

\[
F(x_1,\dots,x_n) =
P(X_1 \le x_1, \dots, X_n \le x_n)
\]

\subsubsection*{Reasoning}

\begin{itemize}
\item Convert point probabilities to cumulative form.
\item Evaluate joint inequalities.
\item Function is non-decreasing in each argument.
\end{itemize}

\subsection{Joint PDF (Continuous Case)}

\[
f(x_1,\dots,x_n) =
\frac{\partial^n}{\partial x_1 \cdots \partial x_n}
F(x_1,\dots,x_n)
\]

\subsubsection*{Reasoning}

\begin{itemize}
\item Start with joint CDF.
\item Differentiate with respect to all variables.
\item Obtain joint density.
\end{itemize}

\subsection{Probability Computation}

\[
P(a_i \le X_i \le b_i) =
\int_{a_1}^{b_1} \cdots \int_{a_n}^{b_n}
f(x_1,\dots,x_n)\, dx_1 \cdots dx_n
\]

\section{Independence of Random Variables}

\subsection{Definition}

Random variables $X_1,\dots,X_n$ are independent if:

\[
p(x_1,\dots,x_n) =
\prod_{i=1}^{n} p_{X_i}(x_i)
\]

or

\[
f(x_1,\dots,x_n) =
\prod_{i=1}^{n} f_{X_i}(x_i)
\]

\subsubsection*{Step-by-Step Reasoning}

\begin{enumerate}
\item Compute joint distribution.
\item Compute marginal distributions.
\item Multiply marginals.
\item Compare:
\begin{itemize}
\item Equal $\Rightarrow$ independent
\item Not equal $\Rightarrow$ dependent
\end{itemize}
\end{enumerate}

\section{Example 2: Uniformly Distributed 3D Vector}

\subsection{Setup}

Let $(X,Y,Z)$ be uniformly distributed in a region.

\subsection{Joint PDF}

\[
f(x,y,z) =
\begin{cases}
c, & \text{inside region} \\
0, & \text{outside}
\end{cases}
\]

\subsection{Derivation}

\begin{itemize}
\item Uniform $\Rightarrow$ constant density $c$
\item Apply normalization:
\end{itemize}

\[
\int_{\text{region}} c\, dV = 1
\]

\[
c = \frac{1}{\text{Volume of region}}
\]

\section{Example 3: Trivariate Random Variables}

\subsection{Joint CDF}

\[
F(x,y,z) = P(X \le x, Y \le y, Z \le z)
\]

\subsection{Joint PDF}

\[
f(x,y,z) =
\frac{\partial^3}{\partial x \partial y \partial z}
F(x,y,z)
\]

\subsection{Marginal Distributions}

\[
f_X(x) = \int \int f(x,y,z)\, dy\, dz
\]

\[
f_Y(y) = \int \int f(x,y,z)\, dx\, dz
\]

\[
f_Z(z) = \int \int f(x,y,z)\, dx\, dy
\]

\subsubsection*{Reasoning}

\begin{itemize}
\item Start with joint PDF.
\item Integrate out unwanted variables.
\item Obtain marginals.
\end{itemize}

\section{Logical Structure of the Lecture}

\begin{itemize}
\item Define joint distribution
\item Introduce PMF, CDF, PDF
\item Study independence
\item Apply to engineering systems (MIMO)
\item Solve examples:
\begin{itemize}
\item General joint case
\item Uniform 3D vector
\item Trivariate variables
\end{itemize}
\end{itemize}

\section{Exam-Focused Key Points}

\begin{itemize}
\item Joint CDF is the foundational definition
\item PMF and PDF derive from it
\item Independence $\Rightarrow$ factorization into marginals
\item Marginals obtained via integration/summation
\item Uniform distributions require normalization
\end{itemize}

\end{document}